\chapter{Performance Measurement of Programs}\label{ch03:performance-measurement-of-programs}

\section{Introduction}

In Chapter 2 we analyzed algorithms theoretically. But theory is not enough. Constants matter. Cache behavior matters. Compiler optimizations matter. The only way to know how fast a program actually runs is to measure it.

This chapter covers practical techniques for measuring program performance using modern tools.

> \textbf{Complete compilable implementations of the data structures in this chapter are in \texttt{code/}.}

\section{Why Measure?}

\begin{itemize}
\item \textbf{Validate theoretical analysis} — does the O(n$^{2}$) algorithm actually scale quadratically?
\item \textbf{Choose between algorithms} — quicksort is O(n log n) on average, but for small arrays insertion sort is faster
\item \textbf{Find bottlenecks} — the 90/10 rule: 90\% of time is spent in 10\% of the code
\item \textbf{Track regressions} — did a change slow things down?
\item \textbf{Tune performance} — iterative improvement requires measurement
\end{itemize}

\section{A First Measurement: std::chrono}

C++11 introduced the \texttt{<chrono>} library for portable high-resolution timing:

\begin{lstlisting}[language=C++]
#include <chrono>
#include <iostream>

class timer {
    using clock = std::chrono::high_resolution_clock;
    using time_point = clock::time_point;
    time_point start_;
public:
    timer() : start_(clock::now()) {}
    void reset() { start_ = clock::now(); }
    double elapsed() const {
        auto end = clock::now();
        return std::chrono::duration<double>(end - start_).count();
    }
};

// Usage:
timer t;
// code to measure
do_something();
std::cout << "Elapsed: " << t.elapsed() << " s\n";
\end{lstlisting}

\subsection{Clock Types}

\begin{center}
\begin{tabular}{cc}
\toprule
Clock & Characteristics \\
\texttt{system\_clock} & System-wide real-time clock, can be adjusted \\
\texttt{steady\_clock} & Monotonic — never adjusted, ideal for intervals \\
\texttt{high\_resolution\_clock} & Highest precision available (often an alias for steady\_clock) \\

\bottomrule
\end{tabular}
\end{center}
\textbf{Always use \texttt{steady\_clock} for performance measurement} — it is immune to system time adjustments.

\subsection{Measuring Short Intervals}

For very short operations, the measurement overhead can dominate. Always measure multiple repetitions:

\begin{lstlisting}[language=C++]
template <typename F>
double measure(F&& f, int iterations = 1000) {
    timer t;
    for (int i = 0; i < iterations; ++i) {
        f();
    }
    return t.elapsed() / iterations;
}
\end{lstlisting}

\section{Google Benchmark Library}

For serious benchmarking, use Google Benchmark. It handles warm-up, iteration count, statistical analysis, and output formatting:

\begin{lstlisting}[language=C++]
#include <benchmark/benchmark.h>

static void BM_VectorPushBack(benchmark::State& state) {
    for (auto _ : state) {
        std::vector<int> v;
        for (int i = 0; i < state.range(0); ++i) {
            v.push_back(i);
        }
    }
}
BENCHMARK(BM_VectorPushBack)->Range(8, 8<<10);

BENCHMARK_MAIN();
\end{lstlisting}

Compile and run:

\begin{lstlisting}
g++ -std=c++20 -O2 -lbenchmark benchmark.cpp -o benchmark
./benchmark
\end{lstlisting}

Output:

\begin{lstlisting}[language=Java]
BM_VectorPushBack/8       25.3 ns  25.1 ns   27692222
BM_VectorPushBack/64       151 ns   148 ns    4716981
BM_VectorPushBack/512     1186 ns  1175 ns     597163
BM_VectorPushBack/4096   10235 ns 10123 ns      68790
BM_VectorPushBack/8192   41234 ns 40987 ns      17021
\end{lstlisting}

\section{Factors Affecting Performance}

\subsection{Hardware Effects}

\begin{itemize}
\item \textbf{CPU speed and architecture} — clock rate, superscalar execution, vector instructions
\item \textbf{Cache hierarchy} — L1 (\textasciitilde{}32KB, \textasciitilde{}1ns), L2 (\textasciitilde{}256KB, \textasciitilde{}4ns), L3 (\textasciitilde{}8MB, \textasciitilde{}10ns), RAM (\textasciitilde{}100ns)
\item \textbf{Memory bandwidth} — sustained throughput to main memory
\item \textbf{Branch prediction} — mispredictions cost 10-20 cycles
\item \textbf{TLB} — virtual address translation cache
\end{itemize}

\subsection{Operating System Effects}

\begin{itemize}
\item \textbf{Context switching} — other processes interrupt your program
\item \textbf{Page faults} — accessing memory not in physical RAM
\item \textbf{CPU frequency scaling} — turbo boost, power saving modes
\end{itemize}

\textbf{Mitigation:} Run benchmarks multiple times, use warm-up phases, run on an idle system.

\subsection{Compiler Effects}

\begin{itemize}
\item \textbf{Optimization level} (\texttt{-O0}, \texttt{-O1}, \texttt{-O2}, \texttt{-O3}, \texttt{-Ofast})
\item \textbf{Inlining} — functions may be inlined or not
\item \textbf{Loop unrolling} — trade code size for speed
\item \textbf{Vectorization} — auto-vectorization of SIMD instructions
\item \textbf{Dead code elimination} — the compiler may remove unused computations
\end{itemize}

\textbf{Important:} Always benchmark with optimizations enabled (\texttt{-O2} or \texttt{-O3}). Debug builds (\texttt{-O0}) give meaningless timing.

\subsection{The Benchmarking Trap: Dead Code Elimination}

The compiler may optimize away the computation you are trying to measure:

\begin{lstlisting}[language=C++]
// WRONG -- compiler may optimize the entire loop away
benchmark::DoNotOptimize(0);  // use this marker instead
static void BM_Bad(benchmark::State& state) {
    int sum = 0;
    for (auto _ : state) {
        for (int i = 0; i < state.range(0); ++i) {
            sum += i;  // compiler may notice 'sum' is unused
        }
    }
}
// The compiler can eliminate the entire inner loop since 'sum' has no observable effect
\end{lstlisting}

Always consume the result:

\begin{lstlisting}[language=C++]
static void BM_Good(benchmark::State& state) {
    int sum = 0;
    for (auto _ : state) {
        for (int i = 0; i < state.range(0); ++i) {
            sum += i;
        }
    }
    benchmark::DoNotOptimize(sum);
}
\end{lstlisting}

\section{Case Study: Sorting Algorithms}

Let us measure insertion sort, merge sort, and std::sort:

\begin{lstlisting}[language=C++]
void insertion_sort(std::span<int> a) {
    for (size_t i = 1; i < a.size(); ++i) {
        int key = a[i];
        int j = static_cast<int>(i) - 1;
        while (j >= 0 && a[j] > key) {
            a[j + 1] = a[j];
            --j;
        }
        a[j + 1] = key;
    }
}

void merge_sort(std::span<int> a) {
    if (a.size() <= 1) return;
    size_t mid = a.size() / 2;
    std::vector<int> left(a.begin(), a.begin() + mid);
    std::vector<int> right(a.begin() + mid, a.end());
    merge_sort(left);
    merge_sort(right);
    std::merge(left.begin(), left.end(), right.begin(), right.end(), a.begin());
}
\end{lstlisting}

\begin{lstlisting}[language=C++]
static void BM_InsertionSort(benchmark::State& state) {
    std::vector<int> v(state.range(0));
    for (auto _ : state) {
        std::iota(v.begin(), v.end(), 0);
        std::shuffle(v.begin(), v.end(), std::mt19937{42});
        insertion_sort(v);
        benchmark::DoNotOptimize(v.data());
    }
}
BENCHMARK(BM_InsertionSort)->Range(8, 8<<10);

static void BM_MergeSort(benchmark::State& state) {
    std::vector<int> v(state.range(0));
    for (auto _ : state) {
        std::iota(v.begin(), v.end(), 0);
        std::shuffle(v.begin(), v.end(), std::mt19937{42});
        merge_sort(v);
        benchmark::DoNotOptimize(v.data());
    }
}
BENCHMARK(BM_MergeSort)->Range(8, 8<<10);

static void BM_StdSort(benchmark::State& state) {
    std::vector<int> v(state.range(0));
    for (auto _ : state) {
        std::iota(v.begin(), v.end(), 0);
        std::shuffle(v.begin(), v.end(), std::mt19937{42});
        std::sort(v.begin(), v.end());
        benchmark::DoNotOptimize(v.data());
    }
}
BENCHMARK(BM_StdSort)->Range(8, 8<<10);
\end{lstlisting}

Expected results:
\begin{itemize}
\item Insertion sort: O(n$^{2}$) — noticeable quadratic curve
\item Merge sort: O(n log n) — smoother curve
\item \texttt{std::sort}: fastest — uses introsort (quicksort + heapsort + insertion sort hybrid)
\end{itemize}

\section{Profiling}

While benchmarking measures total time, \textbf{profiling} shows where time is spent.

\subsection{Sampling Profilers}

Sample the program counter at regular intervals:

\begin{lstlisting}
# macOS
sudo sample ./program -duration 10  # produces a call tree

# Linux
perf record ./program && perf report

# macOS Instruments (GUI)
xcrun xctrace record --template "Time Profiler" --launch ./program
\end{lstlisting}

\subsection{Instrumenting Profilers}

\textbf{Valgrind/Callgrind:}

\begin{lstlisting}
valgrind --tool=callgrind ./program
# Produces callgrind.out.xxx -- view with kcachegrind
\end{lstlisting}

\textbf{Compiler instrumentation (\texttt{-finstrument-functions}):}

GCC/Clang can insert callbacks at every function entry/exit for custom profiling.

\subsection{Cache Analysis with perf}

\begin{lstlisting}
perf stat -e cache-references,cache-misses,cycles,instructions ./program
\end{lstlisting}

Sample output:
\begin{lstlisting}[language=Java]
   1,234,567,890      cycles
     987,654,321      instructions
      12,345,678      cache-references
       1,234,567      cache-misses    # 10.0% miss rate
\end{lstlisting}

A high cache miss rate (>10\%) indicates poor data locality.

\section{Measuring Memory Usage}

\begin{lstlisting}[language=C++]
#include <sys/resource.h>

size_t peak_memory_usage() {
    struct rusage usage;
    getrusage(RUSAGE_SELF, &usage);
    return usage.ru_maxrss;  // in kilobytes (macOS) or bytes (Linux)
}
\end{lstlisting}

For detailed heap profiling, use \textbf{Valgrind/Massif}:

\begin{lstlisting}
valgrind --tool=massif ./program
ms_print massif.out.xxx
\end{lstlisting}

\section{Practical Measurement Methodology}

\begin{enumerate}
\item \textbf{Formulate a hypothesis} — e.g., "std::sort is faster than my quicksort"
\item \textbf{Design the experiment} — what inputs, what sizes, what metrics?
\item \textbf{Warm up} — run the code once to populate caches
\item \textbf{Measure multiple trials} — at least 10, report mean and standard deviation
\item \textbf{Vary input size} — confirm asymptotic class
\item \textbf{Vary input type} — random, sorted, reversed, nearly sorted
\item \textbf{Isolate the system} — close other applications, disable CPU scaling
\item \textbf{Report everything} — compiler, flags, CPU, OS, input parameters
\end{enumerate}

\section{ADT for Polynomial}

An abstract data type for polynomials encapsulates the representation of a polynomial and provides operations to manipulate it. We first specify the ADT, then implement it using a vector of coefficients.

\subsection{ADT Specification}

A polynomial $p(x) = a_{0} + a_{1}x + a_{2}x^{2} + \cdots + a_{n}x^{n}$ is represented by its degree $n$ and a sequence of coefficients $(a_{0}, a_{1}, \ldots, a_{n})$.

\textbf{Operations:}
\begin{itemize}
\item \textbf{Polynomial()} -- construct the zero polynomial
\item \textbf{Polynomial(vector$<$double$>$)} -- construct from a coefficient vector
\item \textbf{degree()} -- return the highest power with a nonzero coefficient
\item \textbf{coefficient(i)} -- return the coefficient of $x^{i}$
\item \textbf{evaluate(x)} -- return $p(x)$ using Horner's method, $O(n)$
\item \textbf{add(q)} -- return $p(x) + q(x)$
\item \textbf{multiply(q)} -- return $p(x) \cdot q(x)$, naive $O(n^{2})$ or FFT-based $O(n \log n)$
\item \textbf{derivative()} -- return $p'(x)$
\item \textbf{operator==} -- test structural equality
\item \textbf{operator$\ll$} -- stream output
\end{itemize}

\textbf{Complexity summary:}
\begin{center}
\begin{tabular}{lc}
\toprule
Operation & Complexity \\
\midrule
evaluate & $O(n)$ \\
add & $O(n)$ \\
multiply & $O(n^{2})$ \\
derivative & $O(n)$ \\
\bottomrule
\end{tabular}
\end{center}

\subsection{Implementation}

We store coefficients in a \texttt{std::vector<double>} where index $i$ holds the coefficient of $x^{i}$. This gives O(1) access to any coefficient and O(n) evaluation.

\begin{lstlisting}[language=C++]
#include <vector>
#include <stdexcept>
#include <iostream>

class polynomial {
public:
    polynomial() = default;

    explicit polynomial(std::vector<double> coeffs)
        : coeffs_(std::move(coeffs)) {
        trim();
    }

    int degree() const {
        return static_cast<int>(coeffs_.size()) - 1;
    }

    bool is_zero() const {
        return coeffs_.empty();
    }

    double coefficient(int i) const {
        if (i < 0 || i > degree()) return 0.0;
        return coeffs_[i];
    }

    // Horner's method: O(n) evaluation
    double evaluate(double x) const {
        if (is_zero()) return 0.0;
        double result = coeffs_[degree()];
        for (int i = degree() - 1; i >= 0; --i) {
            result = result * x + coeffs_[i];
        }
        return result;
    }

    polynomial add(const polynomial& q) const {
        size_t n = std::max(coeffs_.size(), q.coeffs_.size());
        std::vector<double> result(n, 0.0);
        for (size_t i = 0; i < coeffs_.size(); ++i) result[i] += coeffs_[i];
        for (size_t i = 0; i < q.coeffs_.size(); ++i) result[i] += q.coeffs_[i];
        return polynomial(std::move(result));
    }

    // Naive polynomial multiplication: O(n * m)
    polynomial multiply(const polynomial& q) const {
        if (is_zero() || q.is_zero()) return polynomial();
        int n = degree();
        int m = q.degree();
        std::vector<double> result(n + m + 1, 0.0);
        for (int i = 0; i <= n; ++i) {
            for (int j = 0; j <= m; ++j) {
                result[i + j] += coeffs_[i] * q.coeffs_[j];
            }
        }
        return polynomial(std::move(result));
    }

    // Returns the formal derivative: O(n)
    polynomial derivative() const {
        if (degree() <= 0) return polynomial();
        std::vector<double> result(degree());
        for (int i = 1; i <= degree(); ++i) {
            result[i - 1] = i * coeffs_[i];
        }
        return polynomial(std::move(result));
    }

    bool operator==(const polynomial& q) const {
        return coeffs_ == q.coeffs_;
    }

    bool operator!=(const polynomial& q) const {
        return !(*this == q);
    }

    friend std::ostream& operator<<(std::ostream& os, const polynomial& p) {
        if (p.is_zero()) return os << "0";
        for (int i = p.degree(); i >= 0; --i) {
            if (p.coeffs_[i] == 0.0) continue;
            if (i < p.degree()) os << (p.coeffs_[i] > 0 ? " + " : " - ");
            double c = std::abs(p.coeffs_[i]);
            if (c != 1.0 || i == 0) os << c;
            if (i > 0) os << "x";
            if (i > 1) os << "^" << i;
        }
        return os;
    }

private:
    // Remove trailing zero coefficients
    void trim() {
        while (!coeffs_.empty() && coeffs_.back() == 0.0) {
            coeffs_.pop_back();
        }
    }

    std::vector<double> coeffs_;
};
\end{lstlisting}

\subsection{Example Usage}

\begin{lstlisting}[language=C++]
#include <iostream>

int main() {
    // p(x) = 3x^3 + 2x + 1
    polynomial p({1, 2, 0, 3});
    // q(x) = x^2 + 4
    polynomial q({4, 0, 1});

    std::cout << "p(x) = " << p << "\n";
    std::cout << "q(x) = " << q << "\n";
    std::cout << "p(2) = " << p.evaluate(2) << "\n";  // 30

    polynomial sum = p.add(q);
    std::cout << "p+q = " << sum << "\n";

    polynomial prod = p.multiply(q);
    std::cout << "p*q = " << prod << "\n";

    polynomial dp = p.derivative();
    std::cout << "p'(x) = " << dp << "\n";
}
\end{lstlisting}

\subsection{Key Design Decisions}

\begin{itemize}
\item \textbf{Trimming leading zeros} keeps the invariant that the degree is accurate. Every constructor and mutating operation calls \texttt{trim()}.
\item \textbf{Horner's method} for evaluation uses only $n$ multiplications and $n$ additions, with good numerical stability.
\item \textbf{Naive multiplication} is $O(n^{2})$. For high-degree polynomials, the FFT-based approach achieves $O(n \log n)$, but the constant factor is large enough that naive is faster for small degrees.
\item \textbf{Coefficient vector} representation is simple and cache-friendly but wastes space for sparse polynomials (few nonzero terms). A map-based representation would be better for sparse polynomials.
\end{itemize}

\section{ADT for Matrix}

A matrix ADT encapsulates a two-dimensional array of values and provides linear algebra operations. We focus on \textbf{sparse matrices}, which are common in scientific computing and machine learning.

\subsection{ADT Specification}

A matrix $A$ of size $m \times n$ has $m$ rows and $n$ columns. A matrix is \textbf{sparse} when most entries are zero.

\textbf{Operations:}
\begin{itemize}
\item \textbf{Matrix(m, n)} -- construct an $m \times n$ zero matrix
\item \textbf{Matrix(m, n, entries)} -- construct from a list of $(row, col, value)$ triples
\item \textbf{rows(), cols()} -- return dimensions
\item \textbf{get(i, j)} -- return $A_{ij}$, $O(1)$ amortized
\item \textbf{set(i, j, val)} -- set $A_{ij} = \text{val}$
\item \textbf{add(B)} -- return $A + B$, $O(\text{nnz})$
\item \textbf{multiply(B)} -- return $A \cdot B$, $O(\text{nnz}(A) \cdot n_{\text{cols}})$ for naive
\item \textbf{transpose()} -- return $A^{T}$, $O(\text{nnz})$
\end{itemize}

where $\text{nnz}$ is the number of nonzero entries.

\textbf{Complexity summary (sparse representation):}
\begin{center}
\begin{tabular}{lc}
\toprule
Operation & Complexity \\
\midrule
get & $O(\log \text{nnz})$ per row \\
set & $O(\log \text{nnz})$ per row \\
add & $O(\text{nnz})$ \\
multiply & $O(\text{nnz}(A) \cdot n_{\text{cols}}(B))$ \\
transpose & $O(m + n + \text{nnz})$ \\
\bottomrule
\end{tabular}
\end{center}

\subsection{Compressed Row Storage (CRS)}

Compressed Row Storage, also called CSR, is the standard sparse matrix format. It stores:
\begin{itemize}
\item \textbf{values} -- array of all nonzero values, row by row
\item \textbf{col\_index} -- array of column indices corresponding to each value
\item \textbf{row\_ptr} -- array of size $m+1$; \texttt{row\_ptr[i]} is the start of row $i$ in the values array
\end{itemize}

For a $4 \times 5$ matrix:
\[
A = \begin{pmatrix} 1 & 0 & 0 & 2 & 0 \\ 0 & 0 & 3 & 0 & 0 \\ 0 & 4 & 0 & 0 & 5 \\ 0 & 0 & 0 & 6 & 0 \end{pmatrix}
\]

The CRS representation is:
\begin{lstlisting}
values    = { 1, 2, 3, 4, 5, 6 }
col_index = { 0, 3, 2, 1, 4, 3 }
row_ptr   = { 0, 2, 3, 5, 6 }
\end{lstlisting}

Row 2 spans \texttt{values[3..4]} = $\{4, 5\}$ in columns $\{1, 4\}$. Finding a specific element in row $i$ requires a binary search within \texttt{row\_ptr[i]} to \texttt{row\_ptr[i+1]-1}, giving $O(\log \text{nnz\_in\_row})$ access.

\subsection{Implementation}

\begin{lstlisting}[language=C++]
#include <vector>
#include <algorithm>
#include <stdexcept>
#include <iostream>

class sparse_matrix {
public:
    using triplet = std::tuple<int, int, double>;

    sparse_matrix(int rows, int cols)
        : rows_(rows), cols_(cols),
          row_ptr_(rows + 1, 0) {}

    sparse_matrix(int rows, int cols, std::vector<triplet> entries)
        : rows_(rows), cols_(cols) {
        // Sort entries by (row, col)
        std::sort(entries.begin(), entries.end(),
            [](const triplet& a, const triplet& b) {
                if (std::get<0>(a) != std::get<0>(b))
                    return std::get<0>(a) < std::get<0>(b);
                return std::get<1>(a) < std::get<1>(b);
            });

        // Build CRS arrays
        row_ptr_.assign(rows + 1, 0);
        for (const auto& [r, c, v] : entries) {
            if (v != 0.0) {
                row_ptr_[r + 1]++;
                values_.push_back(v);
                col_index_.push_back(c);
            }
        }

        // Convert counts to prefix sums
        for (int i = 0; i < rows; ++i) {
            row_ptr_[i + 1] += row_ptr_[i];
        }
    }

    int rows() const { return rows_; }
    int cols() const { return cols_; }
    int nnz() const { return static_cast<int>(values_.size()); }

    double get(int r, int c) const {
        int start = row_ptr_[r];
        int end = row_ptr_[r + 1];
        auto it = std::lower_bound(
            col_index_.begin() + start,
            col_index_.begin() + end, c);
        if (it != col_index_.begin() + end && *it == c) {
            return values_[it - col_index_.begin()];
        }
        return 0.0;
    }

    // O(nnz) matrix addition
    sparse_matrix add(const sparse_matrix& B) const {
        if (rows_ != B.rows_ || cols_ != B.cols_)
            throw std::invalid_argument("dimension mismatch");

        sparse_matrix result(rows_, cols_);
        std::vector<triplet> entries;

        for (int i = 0; i < rows_; ++i) {
            int a_start = row_ptr_[i], a_end = row_ptr_[i + 1];
            int b_start = B.row_ptr_[i], b_end = B.row_ptr_[i + 1];
            int ai = a_start, bi = b_start;
            while (ai < a_end || bi < b_end) {
                int ac = (ai < a_end) ? col_index_[ai] : cols_;
                int bc = (bi < b_end) ? B.col_index_[bi] : cols_;
                if (ac == bc) {
                    double sum = values_[ai] + B.values_[bi];
                    if (sum != 0.0) entries.emplace_back(i, ac, sum);
                    ++ai; ++bi;
                } else if (ac < bc) {
                    entries.emplace_back(i, ac, values_[ai]);
                    ++ai;
                } else {
                    entries.emplace_back(i, bc, B.values_[bi]);
                    ++bi;
                }
            }
        }
        return sparse_matrix(rows_, cols_, entries);
    }

    // O(nnz(A) * n_cols(B)) sparse matrix multiply
    sparse_matrix multiply(const sparse_matrix& B) const {
        if (cols_ != B.rows_)
            throw std::invalid_argument("dimension mismatch");

        std::vector<triplet> entries;
        for (int i = 0; i < rows_; ++i) {
            std::vector<double> tmp(B.cols_, 0.0);
            for (int ai = row_ptr_[i]; ai < row_ptr_[i + 1]; ++ai) {
                int k = col_index_[ai];
                double a_val = values_[ai];
                for (int bj = B.row_ptr_[k]; bj < B.row_ptr_[k + 1]; ++bj) {
                    tmp[B.col_index_[bj]] += a_val * B.values_[bj];
                }
            }
            for (int j = 0; j < B.cols_; ++j) {
                if (tmp[j] != 0.0) entries.emplace_back(i, j, tmp[j]);
            }
        }
        return sparse_matrix(rows_, B.cols_, entries);
    }

    // O(m + n + nnz) transpose
    sparse_matrix transpose() const {
        sparse_matrix result(cols_, rows_);
        result.row_ptr_.assign(cols_ + 1, 0);

        // Count entries per column (now row in result)
        for (int c : col_index_) {
            result.row_ptr_[c + 1]++;
        }
        for (int i = 0; i < cols_; ++i) {
            result.row_ptr_[i + 1] += result.row_ptr_[i];
        }

        result.values_.resize(nnz());
        result.col_index_.resize(nnz());
        std::vector<int> pos(result.row_ptr_.begin(), result.row_ptr_.end());

        for (int i = 0; i < rows_; ++i) {
            for (int j = row_ptr_[i]; j < row_ptr_[i + 1]; ++j) {
                int c = col_index_[j];
                int dest = pos[c]++;
                result.values_[dest] = values_[j];
                result.col_index_[dest] = i;
            }
        }
        return result;
    }

    friend std::ostream& operator<<(std::ostream& os, const sparse_matrix& m) {
        for (int i = 0; i < m.rows_; ++i) {
            for (int j = 0; j < m.cols_; ++j) {
                if (j > 0) os << " ";
                os << m.get(i, j);
            }
            os << "\n";
        }
        return os;
    }

private:
    int rows_, cols_;
    std::vector<double> values_;
    std::vector<int> col_index_;
    std::vector<int> row_ptr_;
};
\end{lstlisting}

\subsection{When to Use CRS vs Dense}

\begin{center}
\begin{tabular}{lcc}
\toprule
Criterion & Dense (2D vector) & CRS \\
\midrule
Memory & $O(mn)$ & $O(m + n + \text{nnz})$ \\
Access $A_{ij}$ & $O(1)$ & $O(\log \text{nnz\_in\_row})$ \\
Iteration over row & $O(n)$ & $O(\text{nnz\_in\_row})$ \\
Cache behavior & Excellent for small matrices & Good when rows fit in cache \\
Overhead & None & Index arrays \\
\bottomrule
\end{tabular}
\end{center}

Use CRS when the matrix is large and sparse ($\text{nnz} \ll mn$). For matrices that are more than 30-50\% nonzero, a dense representation is usually faster due to better cache locality and no indirection overhead.

\section{ADT for Graph}

A graph $G = (V, E)$ consists of a set of vertices $V$ and a set of edges $E \subseteq V \times V$. The \textbf{Graph Abstract Data Type} specifies the operations available on a graph without committing to a particular representation.

\subsection{ADT Specification}

The Graph ADT supports the following core operations:

\begin{center}
\begin{tabular}{lll}
\toprule
Operation & Signature & Description \\
\midrule
add\_edge & \texttt{add\_edge(u, v)} & Add edge $(u, v)$ \\
remove\_edge & \texttt{remove\_edge(u, v)} & Remove edge $(u, v)$ \\
has\_edge & \texttt{has\_edge(u, v) -> bool} & Test if $(u, v)$ exists \\
degree & \texttt{degree(v) -> int} & Number of neighbors of $v$ \\
bfs\index{BFS} & \texttt{bfs(start) -> vector} & Breadth-first traversal \\
dfs\index{DFS} & \texttt{dfs(start) -> vector} & Depth-first traversal \\
\bottomrule
\end{tabular}
\end{center}

\subsection{Representations}

Two common representations exist:

\begin{itemize}
\item \textbf{Adjacency Matrix.} A 2D array $A$ of size $|V| \times |V|$ where $A[u][v] = 1$ if edge $(u,v)$ exists, 0 otherwise. Space: $O(|V|^2)$. Edge test: $O(1)$. Adding/removing an edge: $O(1)$.
\item \textbf{Adjacency List.} An array of lists, one per vertex, storing neighbors. Space: $O(|V| + |E|)$. Edge test: $O(\deg(u))$. Adding an edge: $O(1)$ (prepend).
\end{itemize}

The adjacency list is preferred for sparse graphs; the adjacency matrix is preferred for dense graphs or when $O(1)$ edge queries are critical.

\subsection{C++ Implementation (Adjacency List)}

\begin{lstlisting}[language=C++, caption={Graph ADT using adjacency list}]
#include <iostream>
#include <vector>
#include <list>
#include <queue>
#include <algorithm>

class Graph {
public:
    explicit Graph(int num_vertices)
        : adj_(num_vertices) {}

    void add_edge(int u, int v) {
        adj_[u].push_back(v);
        adj_[v].push_back(u);
    }

    void remove_edge(int u, int v) {
        adj_[u].remove(v);
        adj_[v].remove(u);
    }

    bool has_edge(int u, int v) const {
        for (int neighbor : adj_[u]) {
            if (neighbor == v) return true;
        }
        return false;
    }

    int degree(int v) const {
        return static_cast<int>(adj_[v].size());
    }

    std::vector<int> bfs(int start) const {
        std::vector<int> visited(adj_.size(), 0);
        std::vector<int> order;
        std::queue<int> q;
        visited[start] = 1;
        q.push(start);
        while (!q.empty()) {
            int v = q.front(); q.pop();
            order.push_back(v);
            for (int w : adj_[v]) {
                if (!visited[w]) {
                    visited[w] = 1;
                    q.push(w);
                }
            }
        }
        return order;
    }

    std::vector<int> dfs(int start) const {
        std::vector<int> visited(adj_.size(), 0);
        std::vector<int> order;
        dfs_helper(start, visited, order);
        return order;
    }

    int num_vertices() const {
        return static_cast<int>(adj_.size());
    }

private:
    std::vector<std::list<int>> adj_;

    void dfs_helper(int v, std::vector<int>& visited,
                    std::vector<int>& order) const {
        visited[v] = 1;
        order.push_back(v);
        for (int w : adj_[v]) {
            if (!visited[w]) {
                dfs_helper(w, visited, order);
            }
        }
    }
};
\end{lstlisting}

\subsection{Complexity Summary}

\begin{center}
\begin{tabular}{lcc}
\toprule
Operation & Adj. List & Adj. Matrix \\
\midrule
Space & $O(|V|+|E|)$ & $O(|V|^2)$ \\
add\_edge & $O(1)$ & $O(1)$ \\
remove\_edge & $O(\deg u)$ & $O(1)$ \\
has\_edge & $O(\deg u)$ & $O(1)$ \\
degree & $O(1)$ & $O(|V|)$ \\
BFS/DFS & $O(|V|+|E|)$ & $O(|V|^2)$ \\
\bottomrule
\end{tabular}
\end{center}

For a graph with $|V| = 10{,}000$ vertices and $|E| = 20{,}000$ edges, the adjacency list uses roughly 240\,KB while the adjacency matrix uses 400\,MB --- a factor of 1{,}700$\times$.

\section{Common Pitfalls}

\begin{center}
\begin{tabular}{ccc}
\toprule
Pitfall & Consequence & Solution \\
Debug build & Measurements include assertions & Use \texttt{-O2} \\
No warm-up & Cold cache skews results & Run iterations before measuring \\
Too few iterations & High variance & Use enough iterations for stable mean \\
Dead code elimination & Zero time reported & Consume results with \texttt{DoNotOptimize} \\
System load & Inconsistent timing & Run on idle system, report median \\
Clock resolution & Zero elapsed time & Use enough work per measurement \\
Inlining & Function not measured separately & Use \texttt{\_\_attribute\_\_((noinline))} \\

\bottomrule
\end{tabular}
\end{center}
\section{Summary}

\begin{enumerate}
\item \textbf{Use \texttt{std::chrono::steady\_clock}} for portable timing.
\item \textbf{Use Google Benchmark} for serious measurement — it handles warm-up, iteration control, and statistics.
\item \textbf{Always optimize} (\texttt{-O2} or \texttt{-O3}) when benchmarking.
\item \textbf{Prevent dead code elimination} by consuming results.
\item \textbf{Profile} to find bottlenecks; \textbf{measure} to verify improvements.
\item \textbf{Understand hardware effects} — cache misses dominate for large data sets.
\item \textbf{Report methodology} so results are reproducible.
\end{enumerate}

\section{Exercises}

\subsection{Drill}

\begin{enumerate}
\item \textbf{Average vs tail latency.} A performance test reports average latency of 2 ms but P99 (99th percentile) of 200 ms. Why are these so different? For a checkout service that must respond within 10 ms, which metric matters? What fraction of customers experience a timeout?
\end{enumerate}

   \textit{(In production, companies like Google and Amazon use P99.9 or P99.99 for SLOs because averages hide the slow requests that hurt user experience.)}

\begin{enumerate}
\item \textbf{Clock choice.} You measure a function with \texttt{std::chrono::steady\_clock} and get 4.8 seconds. A colleague uses \texttt{std::chrono::system\_clock} and gets 5.2 seconds. Both ran the same code on the same machine. What property of \texttt{system\_clock} (vs \texttt{steady\_clock}) explains the difference? Which is correct for benchmarking?
\end{enumerate}

\begin{enumerate}
\item \textbf{Tail latency in servers.} A web server handles 100 concurrent connections. The average response time reported is 2 ms, but P99.9 is 500 ms. An engineer says "the 500 ms is just an outlier — I'll optimize the 2 ms average." Why is this wrong for an e-commerce site? If 0.1\% of requests are slow and the site serves 1 million requests/day, how many users are affected?
\end{enumerate}

\begin{enumerate}
\item \textbf{Dead code elimination.} The following benchmark always reports 0 ns:
\begin{lstlisting}
   static void BM_Sort(benchmark::State& state) {
       std::vector<int> v(1000000);
       std::iota(v.begin(), v.end(), 0);
       for (auto _ : state) {
           std::sort(v.begin(), v.end());  // result never used!
       }
   }
\end{lstlisting}
\end{enumerate}
   Why does the compiler optimize away the sort? How do you fix it?

\begin{enumerate}
\item \textbf{Wall time vs CPU time.} A benchmark shows wall time = 10 s but CPU time = 4 s. Where did the other 6 seconds go? List three possible causes.
\end{enumerate}

   \textit{(In production, burstable cloud instances (like AWS t3) can cause this discrepancy when CPU credits run out.)}

\subsection{Application}

\begin{enumerate}
\item \textbf{Throughput-latency measurement.} A message queue\index{queue} ingests data from producers. On a fast server it handles 800 MB/s. On a modest server, design a benchmark that:
\end{enumerate}
\begin{itemize}
\item Measures throughput (MB/s)
\item Reports P50, P99, P99.9 latency
\item Runs for at least 60 seconds to steady state
\item Compares synchronous (wait for confirmation) vs asynchronous (fire-and-forget) modes
\end{itemize}
   Plot throughput vs latency. What throughput can you guarantee if latency must stay under 50 ms?

\begin{enumerate}
\item \textbf{Load testing.} A web service gets 1000 req/s at night and 10000 req/s during peak. Load testing at peak shows 12 ms average latency. A colleague says "run at 60\% CPU to leave headroom." Another says "90\% CPU is fine, latency is 11 ms." Design an experiment that gradually increases load and finds the inflection point where latency starts to spike.
\end{enumerate}

\begin{enumerate}
\item \textbf{Cache effects.} An analytical query runs 10$\times $ slower the first time than subsequent times. The difference is the cache. Design a benchmark that:
\end{enumerate}
\begin{itemize}
\item Warms the cache (run once, discard)
\item Then measures cold cache (drop OS caches)
\item Tests concurrent queries (1, 4, 16 simultaneous)
\item Tests different data sizes (1 GB, 10 GB)
\item Runs each configuration 10+ times for statistical significance
\end{itemize}
   Report which factor (concurrency, data size, cache state) has the largest impact.

   \textit{(In production, databases like ClickHouse are sensitive to all three — their documentation specifies which knobs to tune.)}

\begin{enumerate}
\item \textbf{SIMD vs scalar.} Compare an element-wise operation on a large array using plain C++ vs using SIMD intrinsics (SSE/AVX on x86 or NEON on ARM):
\end{enumerate}
\begin{itemize}
\item Element-wise multiplication of 1$0^{7}$ floats
\item Dot product of two vectors
\end{itemize}
   Use your platform's timing. Plot time vs array size. At what size does SIMD pull ahead?

\begin{enumerate}
\item \textbf{False sharing.} Two threads increment adjacent variables in an array. Performance is poor beyond 2 cores. The cause is false sharing: the CPU cache line (typically 64 bytes) holds both variables, so each write invalidates the other thread's copy. Design a benchmark that:
\end{enumerate}
\begin{itemize}
\item Creates an array of atomic counters
\item Launches 2, 4, 8 threads, each incrementing its own counter
\item Measures throughput with and without padding (put each counter in its own cache line)
\item Reports speedup vs thread count
\end{itemize}

     \textit{(In production, this pattern appears wherever threads update independent but adjacent data — Uber's H3 geospatial indexing library is one example.)}

\subsection{ADT Design and Implementation}

\begin{enumerate}
\item \textbf{Polynomial evaluation.} Implement the polynomial class from this chapter. Given $p(x) = 5x^{4} - 3x^{2} + 7$, compute $p(1)$, $p(-2)$, and $p(0)$. Verify by hand.
\end{enumerate}

\begin{enumerate}
\item \textbf{Horner vs naive.} Write two evaluation functions: one using Horner's method and one computing each power of $x$ separately using \texttt{std::pow}. Benchmark both for a degree-100 polynomial over $10^{6}$ random inputs. Which is faster and why?
\end{enumerate}

\begin{enumerate}
\item \textbf{Sparse matrix construction.} Create a $5 \times 5$ identity matrix in CRS format. Print the \texttt{values}, \texttt{col\_index}, and \texttt{row\_ptr} arrays. Verify that the number of entries in \texttt{row\_ptr} is $m + 1$.
\end{enumerate}

\begin{enumerate}
\item \textbf{Matrix transpose.} Given a $3 \times 4$ sparse matrix with 5 nonzero entries, compute its transpose by hand. Then run the \texttt{transpose()} method and verify the output matches.
\end{enumerate}

\begin{enumerate}
\item \textbf{Polynomial derivative chain rule.} Extend the polynomial class with a method \texttt{compose(q)} that computes $p(q(x))$. What is the time complexity in terms of the degrees of $p$ and $q$? Implement it and test with $p(x) = x^{2} + 1$ and $q(x) = 2x + 3$.
\end{enumerate}

\begin{enumerate}
\item \textbf{Sparse matrix multiply benchmark.} Generate two random $1000 \times 1000$ sparse matrices with 1\% nonzero density. Measure the time for naive multiply ($O(mn^{2})$ using dense representation) vs CRS multiply. At what density does CRS become faster?
\end{enumerate}

\begin{enumerate}
\item \textbf{Polynomial ADT with map storage.} Implement a second polynomial class that stores nonzero coefficients in a \texttt{std::map<int, double>} instead of a vector. Compare memory usage and evaluation speed with the vector-based version for a degree-1000 polynomial with only 5 nonzero terms.
\end{enumerate}

\begin{enumerate}
\item \textbf{CRS vs COO format.} Implement the Coordinate (COO) sparse matrix format, which stores a list of (row, col, value) triples. Compare it with CRS for: (a) construction time from a stream of entries, (b) element access time, (c) memory usage. When is COO preferable to CRS?
\end{enumerate}

\subsection{Inheritance and Templates}

\begin{enumerate}
\item \textbf{Inheritance hierarchy for matrices.} Design a base class \texttt{matrix\_base} with pure virtual methods \texttt{get}, \texttt{set}, \texttt{rows}, \texttt{cols}. Derive \texttt{dense\_matrix} and \texttt{sparse\_matrix} (using CRS) from it. Write a function \texttt{multiply(matrix\_base\& a, matrix\_base\& b)} that works polymorphically. Discuss the overhead of virtual dispatch in tight loops.
\end{enumerate}

\begin{enumerate}
\item \textbf{Template matrix class.} Write a template class \texttt{matrix<T>} that works for \texttt{int}, \texttt{double}, and a user-defined \texttt{rational} type. Ensure that all arithmetic operations (+, *, transpose) compile and produce correct results for each type.
\end{enumerate}

\begin{enumerate}
\item \textbf{Operator overloading for matrices.} Extend the \texttt{sparse\_matrix} class with operator overloading: \texttt{operator+}, \texttt{operator*}, \texttt{operator==}, and \texttt{operator$\ll$}. Discuss the design tension between returning by value (safe but copies) and returning by reference (efficient but lifetime issues). Implement both and benchmark the difference for large matrices.
\end{enumerate}

\subsection{Research}

\begin{enumerate}
\item \textbf{Cycle-accurate timing.} Research the \texttt{rdtscp} CPU instruction that reads the cycle counter. Write a C++ wrapper. Benchmark \texttt{std::chrono::steady\_clock::now()} overhead vs your \texttt{rdtscp} wrapper over 1$0^{6}$ calls. Which is more precise? When would a cycle counter give misleading results (hint: CPU frequency scaling)?
\end{enumerate}

\begin{enumerate}
\item \textbf{Profiling a real system.} Download an open-source database engine (e.g., DuckDB). Build it. Run a standard benchmark query and profile using \texttt{perf record} or Instruments. Identify:
\end{enumerate}
\begin{itemize}
\item The hottest 3 functions
\item Percentage of time in hash lookups vs parsing vs I/O
\item Cache miss rates (L1, LLC)
\item Branch misprediction rate
\end{itemize}
    Propose one optimization based on the profile.

\begin{enumerate}
\item \textbf{Statistical significance.} Run a sorting benchmark 30 times under identical conditions. Compute mean, std dev, and 95\% confidence interval. Then introduce a small change (e.g., switch from \texttt{std::sort} to \texttt{std::stable\_sort}) — can you detect the difference with 30 runs? With 100 runs? How small a change can you reliably detect?
\end{enumerate}

    \textit{(In production, Google's performance lab uses A/A testing (both groups get the same treatment) to measure the noise floor before running actual experiments.)}

\begin{enumerate}
\item \textbf{Compiler profiling.} Compile a medium-sized C++ project and use your compiler's time-tracing flags (e.g., \texttt{-ftime-report} in GCC/Clang) to measure which optimization pass takes the most time. Is it inlining? Loop optimization? Register allocation? Propose one improvement.
\end{enumerate}

\begin{enumerate}
\item \textbf{FFT-based polynomial multiplication.} Research the Cooley-Tukey FFT algorithm. Implement polynomial multiplication using FFT: (1) convert coefficient vectors to point-value representation via FFT, (2) multiply pointwise, (3) convert back via inverse FFT. Benchmark against naive multiplication for polynomials of degree $2^{10}$ through $2^{16}$. At what degree does FFT become faster?
\end{enumerate}

\begin{enumerate}
\item \textbf{Strassen matrix multiplication.} Research Strassen's algorithm, which multiplies $n \times n$ matrices in $O(n^{2.807})$ time. Implement it for dense matrices. Compare with naive $O(n^{3})$ multiplication for $n = 64, 128, 256, 512$. What is the crossover point? What are the constant-factor penalties?
\end{enumerate}

\begin{enumerate}
\item \textbf{Compressed Column Storage (CCS).} Research CCS, the column-oriented dual of CRS. Implement it. For which operations (row access, column access, transpose, multiply) is CCS faster than CRS? When would a library choose CCS over CRS?
\end{enumerate}

\begin{enumerate}
\item \textbf{Expression templates for lazy evaluation.} Research expression templates, a C++ technique that avoids temporary matrices in expressions like \texttt{C = A + B * D}. Implement a minimal expression template system for matrix operations. Measure the speedup compared to eager evaluation on a chain of 10 matrix additions and multiplications.
\end{enumerate}

\begin{enumerate}
\item \textbf{Polynomial root finding.} Research Newton's method for finding real roots of a polynomial. Implement it using the polynomial class's \texttt{evaluate} and \texttt{derivative} methods. Find all real roots of $p(x) = x^{3} - 6x^{2} + 11x - 6$ (which factors as $(x-1)(x-2)(x-3)$). Discuss convergence guarantees and failure cases.
\end{enumerate}

\section{Primitive and Non-Primitive Data Structures}

\label{sec:primitive-nonprimitive}

Every program manipulates data. The fundamental question is: \emph{how} do we organize that data? Data structures fall into two broad categories: \textbf{primitive} and \textbf{non-primitive}.

\subsection{Primitive Data Structures}

\textbf{Primitive data structures} store a single value directly in memory. They are built into the C++ language and have fixed sizes determined by the hardware and compiler. The CPU can access them in a single operation (typically one memory load).

\begin{lstlisting}[language=C++]
int x = 42;              // 4 bytes, stored in a register or stack
double pi = 3.14159;     // 8 bytes
char c = 'A';            // 1 byte
bool flag = true;        // 1 byte
float f = 2.718f;        // 4 bytes
long long big = 1e18;    // 8 bytes
void* ptr = &x;          // 8 bytes (64-bit pointer)
\end{lstlisting}

\textbf{Properties of primitives:}
\begin{itemize}
\item Fixed size known at compile time
\item Stored by value (no indirection)
\item O(1) access — a single load/store instruction
\item No methods, no encapsulation (just raw bits)
\item Can be combined into \texttt{struct} for composite values
\end{itemize}

\begin{lstlisting}[language=C++]
// Structs compose primitives into a single unit
struct Point {
    double x;  // 8 bytes
    double y;  // 8 bytes
};             // 16 bytes total, contiguous in memory

Point p = {3.0, 4.0};  // stored as 16 consecutive bytes
\end{lstlisting}

\subsection{Non-Primitive Data Structures}

\textbf{Non-primitive data structures} organize multiple data elements into relationships. They are built by the programmer (or provided by libraries) and can grow dynamically. They form the core of this book.

Non-primitive structures are divided into two categories:

\subsubsection{Linear Data Structures}

Elements are arranged in a sequence. Each element has at most one predecessor and one successor.

\begin{center}
\begin{tabular}{lll}
\toprule
Structure & Access & Memory \\
\midrule
Array & O(1) by index & Contiguous \\
Linked List & O(n) by position & Scattered nodes \\
Stack\index{stack} & O(1) top only & Array or linked \\
Queue & O(1) front/rear & Array or linked \\
Hash Table & O(1) average by key & Buckets + nodes \\
\bottomrule
\end{tabular}
\end{center}

\begin{lstlisting}[language=C++]
// Array: contiguous, O(1) index access
int arr[5] = {10, 20, 30, 40, 50};
int val = arr[2];  // O(1) -- direct memory offset

// Linked list: scattered, O(n) traversal
struct Node { int data; Node* next; };
Node* head = new Node{10, new Node{20, new Node{30, nullptr}}};
// To reach the 3rd element: head->next->next->data (O(n))
\end{lstlisting}

\subsubsection{Non-Linear Data Structures}

Elements form relationships beyond simple sequences — trees, graphs, and their variants. Each element can connect to multiple others.

\begin{center}
\begin{tabular}{lll}
\toprule
Structure & Relationship & Example \\
\midrule
Binary Tree & Parent $\to$ 2 children & Expression trees, BSTs \\
Multi-way Tree & Parent $\to$ many children & B-trees, tries \\
Graph & Arbitrary connections & Social networks, road maps \\
Heap & Parent $\leq$ children (partial order) & Priority queues \\
Hash Map & Key $\to$ value & Dictionaries, caches \\
\bottomrule
\end{tabular}
\end{center}

\begin{lstlisting}[language=C++]
// Binary tree: each node has up to 2 children
struct TreeNode {
    int val;
    TreeNode* left;
    TreeNode* right;
};

// Graph: adjacency list
std::vector<std::vector<int>> adj(n);  // adj[u] = neighbors of u
adj[0].push_back(1);  // edge 0 -> 1
adj[0].push_back(2);  // edge 0 -> 2
adj[1].push_back(3);  // edge 1 -> 3
\end{lstlisting}

\subsection{Classification Tree}

\begin{center}
\begin{tabular}{lll}
\toprule
Category & Type & Examples \\
\midrule
Primitive & Integral & \texttt{int}, \texttt{char}, \texttt{bool}, \texttt{long long} \\
& Floating-point & \texttt{float}, \texttt{double} \\
& Pointer & \texttt{T*}, \texttt{std::shared\_ptr<T>} \\
& Composite & \texttt{struct}, \texttt{std::pair} \\
\midrule
Non-Primitive & Linear & Array, Linked List, Stack, Queue \\
& Non-Linear & Tree, Graph, Heap, Hash Table \\
\midrule
Design & Active & \texttt{std::vector}, AVL, Segment Tree \\
& Passive & C-array, POD struct \\
\bottomrule
\end{tabular}
\end{center}

\subsection{Comparison}

\begin{center}
\begin{tabular}{lcccc}
\toprule
Aspect & Primitive & Linear & Non-Linear & Active \\
\midrule
Size & Fixed & Dynamic & Dynamic & Dynamic \\
Access & O(1) & O(1) index / O(n) search & Varies & Varies \\
Memory & Contiguous & Varies & Varies & Varies \\
Operations & None & Insert/Delete & Insert/Delete & Encapsulated \\
Cache & Excellent & Good (array) / Bad (list) & Varies & Varies \\
Overhead & None & Pointers & Pointers + bookkeeping & RAII + invariants \\
\bottomrule
\end{tabular}
\end{center}

\subsection{Interview FAQs}

\begin{enumerate}
\item \textbf{Where does the primitive/non-primitive boundary actually break down in practice?}

The real distinction is ownership semantics. A primitive is a value type --- trivially copyable, no invariants. A non-primitive manages state: it owns memory, enforces invariants, or provides operations that mutate internal structure. But the line blurs constantly. \texttt{std::atomic<int>} is a single value but manages a memory ordering contract. \texttt{std::span} is a view over contiguous memory --- no allocation, but it carries size and bounds. In embedded/HPC, you often \emph{want} non-primitive organization (e.g., SOA layout) with primitive-like semantics (no heap, no indirection). The question isn't ``what category does this fall into?'' but ``who owns the invariants?''
\end{enumerate}

\begin{enumerate}
\item \textbf{Why does the compiler treat \texttt{int[3]} differently from \texttt{int} --- and when does that matter?}

An \texttt{int} is a value type: \texttt{sizeof(int)} is defined by the ABI, it's passed in a register, and assignment copies the value. An \texttt{int[3]} decays to a pointer in most contexts --- it doesn't know its own size, can't be assigned, and can't be returned from functions. This matters when you choose between \texttt{std::array<int,3>} (value type with size awareness, assignable, copyable) vs.\ raw arrays. \texttt{std::array} closes the semantic gap: it's non-primitive in organization but primitive in behavior (trivially copyable, no heap allocation, no invariants).
\end{enumerate}

\begin{enumerate}
\item \textbf{When should you reach for a simple struct/array instead of a library container like \texttt{std::vector} or \texttt{std::map}?}

Use raw/POD when: (1) the size is compile-time known and small (stack allocation avoids heap), (2) you need a trivially-copyable type for IPC/shared memory, (3) you're building a higher-level structure and need full control over layout (arena allocators, GPU buffers). Use library containers when: the size is dynamic, you need exception safety, or you want the compiler to enforce invariants. The real cost of \texttt{std::vector<int>} isn't the data --- it's the heap allocation and indirection. For 4 integers, a \texttt{struct \{ int a,b,c,d; \}} is 3x faster in a hot loop.
\end{enumerate}

\begin{enumerate}
\item \textbf{When does a sorted array outperform a balanced BST, despite both having O(log n) lookup?}

A sorted array has perfect cache locality --- binary search touches O(log n) contiguous elements. A balanced BST (AVL, RB-tree) touches O(log n) nodes scattered across the heap, each causing a cache miss ($\sim$100 cycles). For read-heavy workloads with infrequent inserts, a sorted \texttt{std::vector} with \texttt{std::lower\_bound} often beats \texttt{std::set} by 5--10x in practice. Use a BST when you need O(log n) \emph{insertion} in the middle, or when the collection must stay sorted under mutation. The Linux kernel uses a red-black tree for its scheduler, but Google's btree library uses sorted vectors for exactly this cache-locality reason.
\end{enumerate}

\section{Active and Passive Data Structures}

\label{sec:active-passive}

Data structures can be classified as \emph{active} or \emph{passive} based on how they manage memory and provide operations. Understanding this distinction helps you choose the right design for a given problem.

\subsection{Definitions}

\textbf{Active data structures} encapsulate memory management and provide a set of operations (insert, delete, search) that hide the internal representation. They allocate and deallocate memory internally, enforce invariants (e.g., a BST remains balanced), and expose a high-level API. Examples from this book include \texttt{std::vector}, hash tables, AVL trees, and segment trees.

\textbf{Passive data structures} are plain containers with no built-in operations. They store data in raw memory and leave all management to the caller. Examples include C-style arrays, \texttt{struct} arrays, and POD (Plain Old Data) types.

\begin{center}
\begin{tabular}{lll}
\toprule
Aspect & Active & Passive \\
\midrule
Memory management & Internal (RAII, \texttt{new}/\texttt{delete}) & External (caller) \\
Operations & Encapsulated methods & None (manual) \\
Invariants & Enforced by the structure & Caller's responsibility \\
Typical overhead & Higher (indirection, bookkeeping) & Lower (raw access) \\
Examples & \texttt{std::map}, AVL, heap & \texttt{int[]}, POD struct \\
\bottomrule
\end{tabular}
\end{center}

\subsection{Active Data Structure Example: AVL Tree}

An AVL tree manages its own nodes, performs rotations on insert/delete, and maintains the balance invariant. The caller never touches pointers directly:

\begin{lstlisting}[language=C++]
struct avl_node { int key; int height; avl_node* left; avl_node* right; };

avl_node* insert(avl_node* root, int key) {
    if (!root) return new avl_node{key, 0, nullptr, nullptr};
    if (key < root->key) root->left = insert(root->left, key);
    else if (key > root->key) root->right = insert(root->right, key);
    root->height = 1 + max(height(root->left), height(root->right));
    return rebalance(root);  // rotations enforce invariant
}
\end{lstlisting}

The structure \emph{actively} enforces balance: the caller cannot corrupt the BST property without going through the API.

\subsection{Passive Data Structure Example: C-Array}

A C-style array stores elements contiguously but provides no operations:

\begin{lstlisting}[language=C++]
int arr[100];
int n = 0;

// Caller must manage size, bounds, sorting manually
void push(int val) { arr[n++] = val; }
int pop() { return arr[--n]; }
// No bounds checking, no invariants enforced
\end{lstlisting}

The array is \emph{passive}: it stores bytes but does nothing with them.

\subsection{Hybrid Approach: C++ Containers}

Most C++ STL containers are active: \texttt{std::vector}, \texttt{std::list}, \texttt{std::map}. They combine active memory management with passive storage elements:

\begin{lstlisting}[language=C++]
std::vector<int> v;           // active container
v.push_back(42);              // managed internally

struct point { int x, y; };   // passive element type
std::vector<point> pts;       // active container, passive elements
pts.push_back({3, 4});        // element is a plain struct
\end{lstlisting}

This separation is deliberate: the container handles memory while the element type stays simple and cache-friendly.

\subsection{When to Use Each}

Use \textbf{active} structures when:
\begin{itemize}
\item The collection grows or shrinks dynamically.
\item Invariants (sorted order, balance, uniqueness) must be maintained.
\item Safety and correctness are more important than raw speed.
\end{itemize}

Use \textbf{passive} structures when:
\begin{itemize}
\item The size is fixed or known in advance.
\item You need maximum cache efficiency and minimal overhead.
\item You want full control over memory layout (e.g., SoA vs AoS for GPU compute).
\item Building a higher-level active structure on top (e.g., implementing your own allocator).
\end{itemize}

\subsection{Active vs Passive in Modern C++}

C++17 and C++20 blur the boundary with \texttt{constexpr}, \texttt{std::array}, and \texttt{std::span}:

\begin{lstlisting}[language=C++]
constexpr std::array<int, 5> arr = {3, 1, 4, 1, 5};
// Compile-time fixed array: passive but zero-cost

std::vector<int> v = {3, 1, 4, 1, 5};
std::span<int> view = v;  // passive view over active container
// span provides bounds-checked access without owning memory
\end{lstlisting}

A \texttt{std::span} is a passive view: it references memory managed by an active container. This pattern lets you write generic algorithms that work on both active and passive data:

\begin{lstlisting}[language=C++]
template <typename T>
int sum(std::span<const T> s) {
    int total = 0;
    for (const auto& x : s) total += x;
    return total;
}

// Works with vector, array, C-array, or any contiguous range
\end{lstlisting}

\subsection{Interview FAQs}

\begin{enumerate}
\item \textbf{When building a custom container, what invariants should you enforce internally vs.\ leave to the caller?}

Always enforce internally: memory lifetime (RAII), size tracking, and structural invariants (sorted order, balance). Leave to the caller: value-level invariants (e.g., ``all keys are positive'') and thread safety. The principle is: the container should make \emph{impossible states unrepresentable}. If your stack can underflow, your API design is wrong. The cost of enforcement is indirection and bookkeeping --- use passive structures only when you've measured that the overhead matters and can manage correctness yourself (e.g., arena allocators, SoA layouts for SIMD).

\item \textbf{Can a passive struct accidentally become active through its members?}

Yes --- and this is a common source of hidden costs. A \texttt{struct Node \{ std::vector<int> children; \}} is a passive POD, but constructing it triggers a heap allocation, and copying it deep-copies the vector. Smart pointers are worse: \texttt{struct Obj \{ std::shared\_ptr<Data> d; \}} has an atomic reference count that modifies memory on copy and destruction. This is why hot-path code often uses raw pointers + explicit lifetime management, or arena allocation. The struct's \emph{interface} is passive, but its \emph{behavior} is active --- and that mismatch causes surprises.

\item \textbf{Give a concrete example where replacing \texttt{std::vector} with a raw array improved performance by 10x or more.}

GPU compute: CUDA kernels can't dereference \texttt{std::vector} pointers because the data is on the host. You must copy into a raw \texttt{float*} buffer with a known stride. Similarly, in hot SIMD loops (e.g., dot product over 1M floats), \texttt{std::vector}'s bounds-checking in debug mode and the extra indirection through the control block can cost 2--3x. In lock-free data structures (e.g., LMAX Disruptor), the ring buffer is a raw array specifically because you need predictable memory layout for cache-line padding between producer and consumer indices. The ``active'' wrapper is overhead you're paying on every access.

\item \textbf{Why does \texttt{std::span} exist when we already have iterators and pointers?}

\texttt{std::span} bundles a pointer + size into one value type, making it safe to pass by value (unlike raw pointers where the size is decoupled). Unlike iterators, it works with raw arrays, \texttt{std::vector}, \texttt{std::array}, and C arrays uniformly --- no template bloat. The key insight: \texttt{std::span} lets you write a function \emph{once} and call it with any contiguous range, including stack-allocated arrays. This eliminates the common pattern of needing overloads for \texttt{vector<T>} and \texttt{T*}. In modern C++, \texttt{std::span<const T>} is the idiomatic parameter type for ``read-only contiguous data'' --- it's the non-owning equivalent of \texttt{std::vector<T>}.

\item \textbf{When wrapping a raw buffer in a class, what's the first invariant you should enforce --- and what's the common mistake?}

The first invariant is size tracking: the wrapper must know how many elements are valid, not just how much memory is allocated. The common mistake is adding a \texttt{push\_back} that grows the buffer without invalidating existing iterators/pointers --- this silently creates dangling references. RAII is the second invariant: the destructor must free memory. The transition from passive to active is really a decision about \emph{who is responsible for correctness}. If you wrap but don't enforce invariants, you've added indirection with zero benefit. If you over-enforce (e.g., bounds-checking on every access in a hot loop), you've paid a cost for safety you may not need.
\end{enumerate}

\section{References and Further Reading}

\begin{itemize}
\item Google Benchmark documentation: https://github.com/google/benchmark
\item Chandler Carruth, "Tuning C++: Benchmarks, and CPUs, and Compilers! Oh My!" — CppCon 2015 talk.
\item Agner Fog, "Optimizing Software in C++" — https://www.agner.org/optimize/
\item Ulrich Drepper, "What Every Programmer Should Know About Memory" — https://people.freebsd.org/\textasciitilde{}lstewart/articles/cpumemory.pdf
\item John L. Hennessy and David A. Patterson, \textit{Computer Architecture: A Quantitative Approach}, 6th Edition. Morgan Kaufmann, 2017.
\item David Levinthal, "Performance Analysis Guide for Intel Core i7 and Intel Xeon 5500 Processors" — Intel, 2009.
\item Paul E. McKenney, "Memory Barriers: a Hardware View for Software Hackers" — Linux Technical Conference, 2009.
\end{itemize}

